\section{Introduction}
\label{sec:intro}

Quantum computing developes as a groundbreaking advancement that can tackle classicaly intractable probmels. 
Yet current hardware is limited both in terms of quality and quantity. Different error correcting codes have been proposed over the years that 
enable fault-tolerant quantum computation, yet such codes have been only demonstrated at small scales (Google Quantum AI and Collaborators 2024). Quantum error mitigation (QEM) can close this gap 
and enable utility scale quantum applications. The shared narative amongst QEM methods \cite{RevModPhys.95.045005} is governes between a trade-off between a moderate
resource overhead and a reduction of bias on the algorithms result, nevertheless QEM does not improve the quantum computer. On the other hand, the recent rise of generative models has demonstrated
that neural networks can capture complex patterns, especially in sequential task like text generation. In this work we model quantum processes as an autoregressive task, where noise is accumulated
at every layer. Upon that, a physics-inspired architecture is proposed that builds on the work by 

This work is structured in the following way